\par\noindent\textbf{(1)} 代入 \(p=p_0e^{i(kx-\omega t)}\) 后有 \(-\omega^2+c^2k^2=-i\nu\omega k^2\)，即 \(k^2=\omega^2/(c^2-i\nu\omega)\)。选择 \(\operatorname{Im}k>0\) 的前向支，压力幅随 \(x\) 衰减而非增长。
\par\noindent\textbf{(2)} 若 \(\epsilon=\nu\omega/c^2\ll1\)，展开 \(k=(\omega/c)(1-i\epsilon)^{-1/2}\approx(\omega/c)[1+i\epsilon/2-3\epsilon^2/8+\cdots]\)。故 \(\alpha=\operatorname{Im}k\approx\nu\omega^2/(2c^3)\)，压力幅为 \(e^{-\alpha x}\)。
\par\noindent\textbf{(3)} 在弱耗散下声强与压力振幅平方成正比，\(I(x)\approx I(0)e^{-2\alpha x}\)。因此强度的 \(1/e\) 传播长度 \(\ell_I=1/(2\alpha)\approx c^3/(\nu\omega^2)\)，而压力幅的 \(1/e\) 长度为 \(\ell_p=1/\alpha\)，两者相差2倍。
\par\noindent\textbf{(4)} 当 \(\nu\omega/c^2\) 不再小，必须使用完整复根 \(k=\omega/\sqrt{c^2-i\nu\omega}\) 计算相速度和衰减，不能继续采用 \(\alpha\propto\omega^2\) 的首阶式。高频分子松弛、热传导和非线性也需纳入真实声学耗散。
